% ---------------------------------------------------------------------------------------------------------------
% Geometria de la Vida -- Methods and Concepts manuscript for IEEE Transactions on Ultrasonics (T-USON).
% Version v0.4-tex (27 Nov 2026). INTERNAL DRAFT, NOT SUBMITTED.
%
% Source text: docs/manuscrito/MANUSCRIPT_MC_EN.md (v0.3-en). If they disagree, the results documents cited in the
% text prevail, then the markdown draft.
%
% Template: "TUSON LaTeX Template.zip", IEEE Template Selector (template-selector.ieee.org), path
% Transactions, Journals and Letters -> IEEE Transactions on Ultrasonics -> Original Research -> LaTeX
% (the selector offers no separate Methods and Concepts template). Obtained 27 Nov 2026, 4 670 624 bytes,
% SHA-256 9c614865e91fb2185551e817b1abfb6323d33cfe2cf49a6d323f08b063141804.
% The template files (IEEEtran.cls, TUSON.sty, logos) are NOT redistributed in this repository: download them and
% place IEEEtran.cls and TUSON.sty next to this file before compiling.
%
% Figures: run `python scripts/figure_manuscript.py` from the repository root; it writes docs/figures/*.png.
% ---------------------------------------------------------------------------------------------------------------
\documentclass[journal,twoside]{IEEEtran}
\usepackage[T1]{fontenc}
\usepackage{cite}
\usepackage{amsmath,amssymb,amsfonts}
\usepackage{algorithmic}
\usepackage[]{graphicx}
\usepackage[mathlines,switch]{lineno}
\usepackage{multicol,multirow}
\usepackage{epsfig}
\usepackage{array}
% The template also loads \usepackage[mtbold]{mathtime} and \usepackage{linex}. They are not part of a standard
% free TeX distribution; set \tusonfontstrue when compiling where they are available (for example, IEEE's own setup).
\newif\iftusonfonts
\tusonfontsfalse
\iftusonfonts
  \usepackage[mtbold]{mathtime}
  \usepackage{linex}
\fi
\usepackage{TUSON}

\graphicspath{{../../figures/}}
\newcommand{\repo}[1]{[\texttt{#1}]}   % a results document in the public repository
\newcommand{\TBC}[1]{\textbf{[TO BE COMPLETED by the author: #1]}}

% The template's \GA box fixes the width at 6.25 in. The journal's web guidelines ask for a 660 x 295 px image,
% whose height at that width would exceed the template's 2.5 in limit, so the image is fitted by height instead.
\newcommand{\GAfit}[1]{\vspace{-0.9\baselineskip}\htlne\rightskip=4pt
{\colorbox{mycolor}{\parbox{\linewidth}{\centering\vspace{10pt}\includegraphics[height=2.5in]{#1}\vspace{5pt}}}}\\}

%\linenumbers

\markboth{IEEE Transactions on Ultrasonics}{\TBC{AUTHOR} \lowercase{et al.}: Testing Per-Case Reliability Flags}

\title{Testing Per-Case Reliability Flags for Learned Ultrasound Segmentation: A Preregistered, Sealed-Test Protocol
With Worked Examples in 3-D High-Frequency Mouse Embryo Imaging}
\author{{\TBC{author names and IEEE membership}}
\thanks{Manuscript submitted \TBC{date}. \TBC{funding statement, or ``This work received no external funding.''}}
\thanks{\TBC{affiliation, address and e-mail of each author}. No affiliation is stated until the author provides it.}}

\IEEEaftertitletext{\GAfit{manuscript_graphical_abstract_300dpi.png}
\Abstract{\begin{abstract}
A segmentation network that measures well on average still fails on some cases, and the reference needed to tell
which is rarely available. A per-case reliability flag (green: trust the measurement; yellow or red: do not)
addresses this, but such flags are easily overrated when thresholds are tuned on evaluation data, test sets are
reused, failures are too few to show anything, or negative results go unreported. We describe a protocol for testing
such flags: a tolerance and failure definition fixed in advance; a flag rule derived from the tolerance rather than
fitted to data; a domain guard; calibration gates that must be passed before the test set is opened; frozen
parameters whose cryptographic fingerprint is recorded beforehand and checked by the code; a test set opened once;
three-state verdicts; and evaluability and sample-size rules for asking whether a flag separates correct from failed
measurements. We illustrate each step with preregistered studies on a public dataset of three-dimensional
high-frequency ultrasound volumes of mouse embryos. The gates stopped two flags before they consumed the test set; a
guard withheld green where a certificate calibrated on simulations would have given 540 false greens out of 540; a
stability-based flag passed its gate and, on a sealed test set opened once, failed in 1 of 21 green cases; and a
sample-size simulation exposed a defect in a common rule for declaring no separation, which we correct. The
protocol, code and preregistrations are public.
\hblne
\end{abstract}}
\vspace{1\baselineskip}\vspace*{-1pt}
}

\begin{document}

\maketitle

\Highlights[1]{A preregistered, sealed-test protocol for testing per-case reliability flags of learned ultrasound
segmentation, with evaluability rules for the separation question.}

\Highlights[2]{In 3-D high-frequency mouse embryo ultrasound, gates stopped two flags before the test; a stability
flag tested once failed in 1 of 21 green cases (95\% CI 0--14\%).}

\Highlights[3]{Simulation shows a common no-separation rule misfires when failures are few; requiring at least 10
failures guards small validation studies against false conclusions.}

\Keywords{Deep learning, high-frequency ultrasound, mouse embryo, preregistration, quality control,
reproducibility, segmentation, uncertainty estimation}

\PrintHighlights

% =================================================================================================================
\section{Introduction}
\IEEEPARstart{V}{olume} and shape measurements from ultrasound increasingly come from learned segmentation. Even
accurate networks fail on some inputs, and the user rarely has a reference to check a given case. A per-case
reliability output is therefore useful: it tells the user when to trust a number.

Disagreement between predictions on transformed inputs is an established estimate of segmentation uncertainty
\cite{wang2019a}, and it has been used to indicate potential mis-segmentations \cite{wang2019b}. Evaluating whether
such an estimate can support a \emph{decision} about a single case is harder than showing that it correlates with
error. Common pitfalls are:
\begin{itemize}
\item thresholds chosen while looking at the evaluation data;
\item the same data used to design and to judge the flag;
\item test sets reused after a first look;
\item conclusions about whether the flag separates correct from failed cases drawn from a handful of failures;
\item failed attempts left out of the report.
\end{itemize}

This paper documents a protocol that guards against those pitfalls, as a reusable practice rather than a single
result. Section~\ref{sec:protocol} describes the protocol, Section~\ref{sec:examples} works through each step with
real studies, and Section~\ref{sec:discussion} discusses what the protocol costs and what it does not solve.

The examples use a public dataset of 3-D high-frequency ultrasound volumes of mouse embryos with manual
segmentations of the body and brain ventricles \cite{qiu2021,dataport,xu2020}. No new animal data were generated and
no human data were used.

% =================================================================================================================
\section{The Protocol}\label{sec:protocol}

\subsection{Before Any Data Are Seen}\label{sec:before}
\begin{enumerate}
\item \textbf{Fix the tolerance and the failure definition.} Example: a volume measurement fails if it differs from
the reference by more than $\tau = 10\%$. For a small structure, a relative tolerance can be unreachable, and a
tolerance in boundary voxels, $e_b = |V_M - V_R| / A_R \le 0.5$ voxels, is more appropriate
(Section~\ref{sec:exD}).
\item \textbf{Derive the flag rule from the tolerance, not from data.} Example (Section~\ref{sec:exC}): green if the
spread of the measured volume across the 8 flips of the image is at most $\tau/2$ and the disagreement between the
two network stages is at most $\tau$; red if either exceeds $2\tau$ or $3\tau$, respectively; yellow otherwise. With
few failures, any fitted threshold learns noise.
\item \textbf{Add a domain guard.} It is an envelope of simple image indices (here, an estimated signal-to-noise
ratio and the edge width, 1--99\% percentiles) outside which the flag never gives green. Check that the guard is
\emph{computable} on the target structure (Section~\ref{sec:exD}).
\item \textbf{Split by subject} (embryo), never by image, into a part for fitting the network, a calibration part
and a test part.
\item \textbf{Write the preregistration:} hypotheses and criteria; minimum numbers for each hypothesis to be
evaluable; predictions with probabilities; and what each outcome will mean. Upload it with the code and its tests
before computing anything that decides.
\end{enumerate}

\subsection{Gates, Freezing and the Single Test}
\begin{enumerate}\setcounter{enumi}{5}
\item \textbf{Feasibility gate (G0)} on a few fitting cases: the signals must be computable and the pipeline must
reproduce the measurement exactly.
\item \textbf{Calibration gate} on held-out calibration subjects. Example: at least 10 greens, and at most 10\% of
them failing. If the gate fails, \emph{the test set is not opened}, and a new preregistration is needed.
\item \textbf{Freeze and fingerprint.} Write rules, envelope and network hashes to one file, and record its SHA-256
in the preregistration before the test. The test command checks the fingerprint, the network hashes and the code
thresholds, and refuses to run if any differ.
\item \textbf{Open the test once.} Afterwards the test set is \emph{used}: no later study may decide with it.
\item \textbf{Report three states.} A hypothesis is supported, not supported, or not evaluable or inconclusive. A
value that cannot be computed is never ``not supported.''
\end{enumerate}

\subsection{The Separation Question}\label{sec:separation}
Showing that greens rarely fail is not the same as showing that the flag adds information. With a low base failure
rate, almost any rule that leaves enough greens meets a ``greens rarely fail'' criterion. The separation question
compares failure rates between non-greens and greens: $\Delta$ is the failure rate among non-greens minus that among
greens, with a 95\% confidence interval (CI) by percentile bootstrap, resampling subjects. The flag ``separates'' if
the lower bound is above 0, ``does not separate'' if the upper bound is at most 5 points, and the result is
``inconclusive'' otherwise.

We add three rules learned in Section~\ref{sec:exE}:
\begin{itemize}
\item \textbf{S-1, evaluability by failures:} require at least 10 failures in total and at least 1 in each group.
Otherwise the separation hypothesis is not evaluable.
\item \textbf{S-2, planning:} declare the expected failure rate and aim for about 25 expected failures. With fewer,
state in advance that ``inconclusive'' is the likely outcome.
\item \textbf{S-3:} keep resampling by subject.
\end{itemize}

\subsection{Reporting}
Compare the written predictions with the outcomes. Record every deviation with its date, approved before the
affected quantity is computed. Publish negative and non-evaluable studies with the same detail as positive ones.
Table~\ref{tab:checklist} lists the protocol as a checklist.

\begin{table*}[!t]
\centering
\caption{Checklist for Testing a Per-Case Reliability Flag}\label{tab:checklist}
\footnotesize
\renewcommand{\arraystretch}{1.15}
\begin{tabular}{@{}l p{3.3cm} p{3.1cm} p{6.4cm} p{2.6cm}@{}}
\hline
\# & Item & When & Done if & Example in Section~\ref{sec:examples} \\
\hline
1 & Tolerance and failure definition & Before any data & Written in the preregistration, in units the structure can
meet & D (relative $\to$ boundary voxels) \\
2 & Flag rule & Before any data & Thresholds derived from the tolerance; nothing fitted & C \\
3 & Domain guard & Before any data & Indices defined and checked to be computable on the target structure & A, D \\
4 & Split by subject & Before any data & No subject in two parts & All \\
5 & Preregistration & Before computing anything decisive & Hypotheses, criteria, minimums, predictions and meanings
uploaded with code and tests & All \\
6 & Feasibility gate (G0) & On fitting cases & Signals computable; pipeline reproduces the measurement exactly & C \\
7 & Calibration gate & On calibration subjects & Passed; otherwise the test is not opened & B, D \\
8 & Freeze and fingerprint & Before the test & SHA-256 recorded in the preregistration; code refuses mismatches & C,
D \\
9 & Single test & Once & Test opened once, then marked as used & C, D \\
10 & Three-state verdicts & At closing & Supported / not supported / not evaluable or inconclusive; uncomputable is
never ``not supported'' & All \\
S-1 & Separation evaluability & At planning & $\ge 10$ failures in total and $\ge 1$ in each group & E \\
S-2 & Separation planning & At planning & Expected failure rate and number of cases declared; about 25 expected
failures & E \\
S-3 & Resampling unit & At planning & Bootstrap by subject & C, D \\
R & Reporting & At closing & Predictions compared with outcomes; dated deviations; negative results published & All \\
\hline
\end{tabular}
\end{table*}

% =================================================================================================================
\section{Worked Examples}\label{sec:examples}
All studies below follow Section~\ref{sec:protocol}. Each cites its results document in the public repository
\cite{repo}.

\subsection{A Domain Guard Refuses Where a Simulation-Calibrated Certificate Fails}\label{sec:exA}
A certificate calibrated on synthetic data was applied to 1188 real measurements: 180 body and 216 ventricle cases,
each with 3 threshold segmenters. The guard gave no green. Without it, all 540 greens would have been false
\repo{docs/h11/RESULTADOS\_H11.md}. Recalibrating on simulated ultrasound did not change this: 180 of 180 false
greens without the guard \repo{docs/h11/RESULTADOS\_H12R.md}.

\emph{Lesson (step 3):} the guard is what made the certificate safe on a domain it had never seen.

\subsection{A Calibration Gate Stops a Flag Before It Consumes the Test Set}\label{sec:exB}
A two-stage network, reimplemented from \cite{qiu2021,xu2020}, put 24 of 27 calibration cases within 10\% (median
Dice 0.938). A certificate based on image observables then failed the calibration gate: 2 of 16 greens failed, more
than giving green to every case. The test set was not opened \repo{docs/h19/RESULTADOS\_H19.md}.

\emph{Lesson (step 7):} the gate saved the test set for a better candidate.

\subsection{A Stability Flag Passes and Is Tested Once}\label{sec:exC}
The flag of step 2 passed the calibration gate with 12 greens and no failures. After freezing (fingerprint recorded
in the preregistration), the sealed test set (44 cases from 43 embryos) was measured once
\repo{docs/h20/RESULTADOS\_H20.md}; Table~\ref{tab:h20} and Figs.~\ref{fig:lights}(a) and \ref{fig:signals}
summarize the result.

\begin{table}[!t]
\centering
\caption{Sealed Body Test (44 Cases, Opened Once)}\label{tab:h20}
\footnotesize
\begin{tabular}{@{}p{4.0cm} p{4.1cm}@{}}
\hline
Measure & Value \\
\hline
Flags & 21 green, 22 yellow, 1 red \\
Failures among greens (primary) & 1 of 21: 4.8\% (CI 0.0--14.3\%); supported \\
Failures among non-greens & 2 of 23: 8.7\% \\
Difference (non-green $-$ green) & 3.9 points (CI $-11.4$ to 20.0) \\
Failure rate in all cases & 3 of 44: 6.8\% \\
Cases within 10\% & 93.2\% (median Dice 0.939) \\
\hline
\end{tabular}
\end{table}

\emph{Lesson (Section~\ref{sec:separation}):} the primary criterion was met, but with a 6.8\% base failure rate
almost any rule would have met it. The separation interval includes 0. The failing green underestimates volume by
12.9\% with both signals low: the network made the same error on all 8 inputs, which a stability flag cannot see.

\begin{figure*}[!t]
\centering
\includegraphics[width=\textwidth]{manuscript_fig1_lights.png}
\caption{Error of each test case by flag. (a) Body, sealed test (H20, 44 cases): absolute volume error, tolerance
10\%. (b) Brain ventricles, test (H25, 40 cases): boundary error $e_b$, tolerance 0.5 voxels. Crosses are failures.
Generated from already-measured cases by \texttt{scripts/figure\_manuscript.py}, with no new analysis.}
\label{fig:lights}
\end{figure*}

\begin{figure}[!t]
\centering
\includegraphics[width=\columnwidth]{manuscript_fig2_signals.png}
\caption{The two stability signals on the sealed body test set (H20), with the thresholds of the flag rule
($\tau/2$ and $2\tau$ for the flip spread; $\tau$ for the stage disagreement). Crosses are failures
($|\text{error}| > 10\%$).}
\label{fig:signals}
\end{figure}

\subsection{Changing the Structure Exposes the Units of the Tolerance and a Guard Defect}\label{sec:exD}
The same rule, unchanged, gave no green at all on brain ventricles: 0 of 34 in calibration
\repo{docs/h23/RESULTADOS\_H23.md}.

With the tolerance and signals re-expressed in boundary voxels, the calibration gate still failed: 3 greens of 34.
The guard admitted only 5 of 34 cases, because its indices were not computable for a structure darker than the
background (27 of 34 cases) \repo{docs/h24/RESULTADOS\_H24.md}.

A new preregistration corrected the guard to use the absolute contrast. It then admitted 90\% of test cases, against
22.5\% for the old guard. The gates passed, and the ventricle test set was opened once (40 cases)
\repo{docs/h25/RESULTADOS\_H25.md}; see Fig.~\ref{fig:lights}(b). The network was within 0.5 voxels in 95\% of cases;
the primary hypothesis was not evaluable (19 greens, where 20 were required); and separation was inconclusive (1 of
19 greens and 1 of 21 non-greens failed).

\emph{Lessons (steps 1, 3 and 7):} a tolerance must be expressed in units the structure can meet; the guard must be
checked for computability on each new structure; and the gates caught both problems without touching the test set.

\subsection{A Sample-Size Simulation Exposes a Defect in the ``No Separation'' Rule}\label{sec:exE}
To plan a future study, we simulated the separation rule of Section~\ref{sec:separation}
\repo{docs/h25/TAMANO\_MUESTRAL\_SEPARACION.md}.

\emph{Power.} With about 80 cases and a failure gap like that seen in calibration (5\% versus 27\%), the rule detects
separation 83\% of the time. A gap of 10 points needs about 270 cases. What matters is the number of failures, not
the number of cases.

\emph{A defect.} With no true separation and 40 cases, the rule declared ``does not separate'' in 35\% of
simulations. In 140 of those 144 simulations there were no failures among non-greens, and the bootstrap interval
degenerated.

No published verdict in the series was affected, but a future study could have concluded ``no separation'' from this
artefact. Rules S-1 to S-3 (Section~\ref{sec:separation}) prevent it.

% =================================================================================================================
\section{Discussion}\label{sec:discussion}

\subsection{What the Protocol Buys}
Every claim can be traced to a preregistration fixed before the decisive computation. No threshold was changed after
seeing test data. Each sealed test set was opened at most once. Two candidate flags were stopped at their gates
(Sections~\ref{sec:exB} and \ref{sec:exD}) before they could consume a test set.

\subsection{What It Costs}
Test sets are used up: in this series, both the body and the ventricle test sets are now spent. Many studies end as
``not evaluable''; half of the studies on real data in this series did. This is the protocol declining to draw
conclusions the data could not support, not a failure of the protocol.

\subsection{What It Does Not Solve}
A flag based on stability cannot see a consistent error. A flag cannot be shown to separate correct from failed
measurements without enough failures, which an accurate network rarely produces on one dataset. The examples come
from one public dataset, one system and one species; a direct replication with new volumes is preregistered as a
base document \repo{docs/h26/PREREGISTRO\_H26\_BORRADOR.md}.

\subsection{Adapting the Protocol}
Steps 1--10 and rules S-1 to S-3 do not depend on the modality or the network. To reuse them, an author needs a
tolerance meaningful for the structure, signals available without the reference, a guard computable on the target
images, and enough subjects for gates and a single test.

None of the examples concerns humans or clinical use.

% =================================================================================================================
\section{Conclusion}
Per-case reliability flags for learned ultrasound segmentation can be tested without overrating them if the
tolerance, rule and minimums are fixed in advance, gates protect the test set, the frozen instrument is fingerprinted
and the test is opened once. Whether a flag separates correct from failed measurements is a separate question that
needs enough failures to be evaluable; rules S-1 to S-3 make that requirement explicit before the data are collected.

% =================================================================================================================
\section*{Data and Code Availability}
The data are public on IEEE DataPort \cite{dataport} (CC BY) and are not redistributed. Code, preregistrations,
results and the protocol rules (\texttt{docs/REGLAS\_PREREGISTRO.md}) are in the project repository \cite{repo}
(MIT license); reproduction steps are in \texttt{docs/REPRODUCCION.md}.

\section*{Ethics Statement}
Only previously published animal data were used; no new animal experiments were performed and no human data were
used.

\section*{Acknowledgment}
The code and texts were developed with AI assistance (Claude, Anthropic). The work has not had external human
scientific review. \TBC{other acknowledgments}

\section*{Competing Interests}
\TBC{competing interests statement}

% =================================================================================================================
\begin{thebibliography}{6}

\bibitem{qiu2021}
Z.~Qiu \emph{et al.}, ``A deep learning approach for segmentation, classification, and visualization of 3-D
high-frequency ultrasound images of mouse embryos,'' \emph{IEEE Trans. Ultrason., Ferroelectr., Freq. Control},
vol.~68, no.~7, pp.~2460--2471, Jul. 2021, doi: 10.1109/TUFFC.2021.3068156.

\bibitem{dataport}
``Ultrasound mouse embryo segmentation and classification,'' IEEE DataPort. [Online]. Available:
https://doi.org/10.21227/drmd-xq36

\bibitem{xu2020}
T.~Xu \emph{et al.}, ``Deep Mouse: An end-to-end auto-context refinement framework for brain ventricle \& body
segmentation in embryonic mice ultrasound volumes,'' in \emph{Proc. IEEE 17th Int. Symp. Biomed. Imag. (ISBI)},
2020, pp.~122--126, doi: 10.1109/ISBI45749.2020.9098387.

\bibitem{wang2019a}
G.~Wang, W.~Li, M.~Aertsen, J.~Deprest, S.~Ourselin, and T.~Vercauteren, ``Aleatoric uncertainty estimation with
test-time augmentation for medical image segmentation with convolutional neural networks,''
\emph{Neurocomputing}, vol.~335, pp.~34--45, Mar. 2019, doi: 10.1016/j.neucom.2019.01.103.

\bibitem{wang2019b}
G.~Wang, W.~Li, S.~Ourselin, and T.~Vercauteren, ``Automatic brain tumor segmentation based on cascaded
convolutional neural networks with uncertainty estimation,'' \emph{Front. Comput. Neurosci.}, vol.~13, 2019,
Art. no.~56, doi: 10.3389/fncom.2019.00056.

\bibitem{repo}
Geometr\'{\i}a de la Vida project repository, branch \texttt{geovida/v0.1-desde-cero}. [Online]. Available:
\TBC{public repository URL, preferably with an archived release DOI}

\end{thebibliography}

% =================================================================================================================
\begin{IEEEbiographynophoto}{\TBC{Author name}}
\TBC{biography; the journal asks for a biography and photograph of each author, which may be supplied at the end}
\end{IEEEbiographynophoto}

\vfill
\end{document}
